A multi-task leaderboard gives a developer a reason to train a chosen model on data related to some tasks. For two fixed public score tables we ask how much targeted investment makes a target model first, and how this depends on the aggregation rule. The target raises its task scores within a per-task cap while rivals stay fixed; we solve exactly for the fewest tasks touched ($K$), the fewest score points ($L_1$) and the smallest largest single-task rise ($L_\infty$) under five score rules, plus Borda, plurality and a Condorcet requirement on HELM and a top-500 frame. The tables are HELM MMLU (79 models, 57 subjects) and the Open LLM Leaderboard (4{,}496 models, 6 tasks). We prove that no universal ordering of the mean, median and minimum exists. (i) With unrestricted headroom the mean needs the most raw score points but not always the most tasks, and not the most points in standard-deviation units. (ii) With at most $+10$ points per task, the mean can put 29 of 78 HELM targets first (37\%, resampling range 26--49\%; 14\% at $+5$, 72\% at $+20$) and 113 of 4{,}495 Open LLM Leaderboard targets (2.5\%, 2.0--3.0\%; 0.3\% and 18\%). Whether the minimum rule is more reachable than the mean at $+10$ differs between the tables, but not only because of the table: matching rivals or tasks keeps each sign in most random subsets, while a fixed mean-selected frame and a task-IQR cap reverse it. We name no single cause. (iii) Lowest-score-first greedy already attains the optimal number of tasks at the ceiling; random order needs 1.8, 1.1 and 4.2 times the optimum on HELM under the mean, median and minimum. All numbers describe fixed snapshots, not what training achieves.
